\documentclass[11pt,a4paper,UTF8,fontset=fandol]{ctexart}
\usepackage[margin=2.35cm]{geometry}
\usepackage{amsmath,amssymb,mathtools,bm,booktabs,array,graphicx,hyperref,xcolor,enumitem,fancyhdr}
\usepackage{tikz}
\usepackage[section]{placeins}
\usepackage{float}
\AddToHook{cmd/section/before}{\clearpage}
\usetikzlibrary{arrows.meta,positioning,calc}
\tikzset{gene/.style={draw,rounded corners,minimum width=12mm,minimum height=8mm,fill=blue!6},act/.style={-{Stealth},thick,blue!65!black},rep/.style={-{Bar[width=6pt]},thick,red!70!black}}
\hypersetup{colorlinks=true,linkcolor=blue!50!black,urlcolor=blue!45!black,citecolor=blue!45!black}
\setlength{\headheight}{15pt}
\setlength{\parindent}{2em}\setlength{\parskip}{.25em}
\setcounter{tocdepth}{2}\numberwithin{equation}{section}
\newcommand{\dd}{\mathrm d}\newcommand{\R}{\mathbb R}\newcommand{\ssv}{\mathrm{ss}}
\pagestyle{fancy}\fancyhf{}\fancyhead[L]{最小基因调控回路：从反应到动力学}\fancyhead[R]{\thepage}
\title{\bfseries 最小基因调控回路的理论基础\\[4pt]\large 从化学反应到一、二、三组分系统及真实网络}
\author{可复现学习笔记（含源代码与数值图）}
\date{2026 年 10 月}
\begin{document}\maketitle
\begin{abstract}
本文采用明确的反应规则与动力学假设，从产生--降解、启动子结合和质量作用定律出发，研究少组分基因调控回路所允许的行为。全文围绕一个问题展开：调控如何把简单的浓度恢复变成记忆、振荡或适应？先从反应建立表达模型，再依次加入自反馈、相互反馈和第三个调控环节，沿着每个模型的稳态与扰动响应解释这些行为的来源。双节点部分完整整理 81 种符号拓扑，进一步结合 Science（2008）的正负反馈振荡研究与 Cell（2009）的适应网络枚举，比较可调时钟、时间筛选、缓冲反馈与比例前馈。随后把连续浓度改为离散分子数，最后用乳糖操纵子讨论最小回路如何连接真实网络。重点在于区分\emph{结构所允许的机制}、\emph{特定参数下实现的动力学}与\emph{具体生物学功能}，而非追求复杂网络的完整重构。
\end{abstract}
\newpage{\footnotesize\linespread{.93}\selectfont\setlength{\parskip}{0pt}\tableofcontents}\newpage
\section{从一个问题出发：调控怎样改变表达的时间行为}
一个以恒定速率产生、以一阶速率降解的蛋白质，总会逐渐回到同一个浓度。基因调控的作用，是让产生率依赖其他分子，甚至依赖自身：同样的外部条件下，系统于是可能保留不同表达状态，也可能持续振荡，或者只对输入变化作出短暂响应。要理解这些现象，需要依次回答三个问题：产生率从什么反应机制来？反馈如何改变稳态和扰动的演化？这些局部行为在随机性与网络背景中还能否成立？

每个模型都按“回路与问题、方程与推导、适用条件”的顺序讲解。随着动态变量增加，比较新环节如何带来状态选择、振荡或适应，再回到分子计数和真实网络。

基因调控网络的节点常取为基因表达产物（例如蛋白质）的有效浓度，并不意味着微观反应中只含这些分子。启动子状态、RNA 聚合酶、mRNA、配体、代谢物均可能在简化模型时被省略。本文用 $x,y,z$ 表示按参考浓度换算后的数值，例如 $x=1$ 表示达到参考浓度。激活边表示：调控因子增加时，目标的产生率上升；抑制边表示产生率下降。用公式说，就是 $\partial F_i/\partial x_j$ 分别为正或负。边的符号说明作用方向，具体结合方式还要另外建模。

\begin{figure}[H]\centering
\begin{tikzpicture}[node distance=22mm]
\node[gene] (x) {$X$}; \node[gene,right=of x] (y) {$Y$};
\draw[act] (x)--node[above]{激活：$+$}(y);
\node[gene,right=35mm of y] (u) {$X$}; \node[gene,right=of u] (v) {$Y$};
\draw[rep] (u)--node[above]{抑制：$-$}(v);
\end{tikzpicture}
\caption{调控边的记号。箭头表示产生率随调控变量增加，钝头表示产生率下降；边的符号不能单独指定结合机制。}\label{fig:signs}
\end{figure}

标准连续模型为
\begin{equation}\dot x_i=F_i(\bm{x};\bm{u})-\gamma_i x_i,\qquad F_i\geq0,\quad\gamma_i>0.\end{equation}
其中 $F_i$ 为有效产生率，$\bm u$ 为外部输入，$\gamma_i$ 汇总一阶降解及相应的稀释效应。讨论稳态时固定输入；讨论响应时再指定输入怎样变化。同一组初始浓度能否决定唯一一条轨迹，是后面所有相图分析的前提。一个常用的保证条件叫\emph{局部 Lipschitz 连续}：在每个所考察状态附近，变化率 $\bm v(\bm x)$ 满足
\[
\|\bm v(\bm x)-\bm v(\bm y)\|\leq L\|\bm x-\bm y\|,
\]
其中 $L$ 是这个局部范围内的某个有限常数，$\|\cdot\|$ 可以理解为两个状态之间的距离。直观地说，浓度只改动一点，变化率也只能改动有限的一点，不能在很小的浓度差上出现无限陡的变化。这个条件比“曲线不断开”更强，却不要求每个位置都有斜率；它保证同一起点不会分叉出两种未来。例如 $v(x)=|x|$ 满足 $|v(x)-v(y)|\leq|x-y|$，可取 $L=1$；虽然零点处有尖角、没有导数，它仍然是 Lipschitz 连续的。“局部”表示不同浓度范围可以用不同的 $L$，并不是要求整个浓度空间共用一个限制。

使用导数分析小扰动时，本文通常进一步假设产生率为 $C^1$。这个记号表示\emph{函数可求一阶导数，且一阶导数连续}：不仅曲线不断开，斜率本身也连续变化。$C^1$ 函数在足够小的范围内满足上述 Lipschitz 条件，因此本文的 Hill 等常见模型可以同时支持唯一轨迹和斜率分析。这些是假设在所分析的浓度范围内成立；若模型有开关式跳变或边界处无限陡的斜率，就要重新检查。

\emph{稳态}指产生与损失平衡、浓度不再变化。\emph{稳定}指小偏差下轨迹始终留在附近；若偏差逐渐消失，则称为\emph{渐近稳定}。局部结论只涉及稳态附近，不代表所有初始浓度都回到它。



\section{建立共同模型：从反应到有效产生率}
所有回路共享同一个骨架：分子通过反应产生和消耗。先把这个骨架写清楚，再讨论哪些微观变量可以消去，以及调控如何进入产生率。
\subsection{产生与降解给出表达的基线}
对于 $\sum_i\alpha_{ir}X_i\xrightarrow{k_r}\sum_i\beta_{ir}X_i$，质量作用律给出 $v_r(\bm x)=k_r\prod_i x_i^{\alpha_{ir}}$，以及
\begin{equation}\dot{\bm x}=S\bm v(\bm x),\qquad S_{ir}=\beta_{ir}-\alpha_{ir}.\end{equation}
其中 $S$ 记录每次反应使各类分子增加或减少多少，$v_r$ 则表示第 $r$ 个反应发生得多快。例如一次反应消耗一个 $X$、产生两个 $Y$，对应的变化量就是 $-1$ 和 $+2$。对于 $\varnothing\xrightarrow{\alpha}X$, $X\xrightarrow{\gamma}\varnothing$，
\begin{equation}\dot x=\alpha-\gamma x,\qquad x(t)=\frac\alpha\gamma+\left(x(0)-\frac\alpha\gamma\right)e^{-\gamma t}.\end{equation}
这个解把后文的基线完整地展示出来：唯一稳态是 $x^*=\alpha/\gamma$，任何非负初值与它的差都以 $e^{-\gamma t}$ 衰减，恢复时间为 $1/\gamma$。没有反馈时，系统不会选择多个长期状态。

\subsection{把 mRNA 隐去，究竟省略了什么}
蛋白质的产生通常还经过转录和翻译。保留 mRNA 时，反应与方程为
\begin{align}
\varnothing\xrightarrow{k_m}M,\quad M\xrightarrow{\delta_m}\varnothing,\quad M\xrightarrow{k_p}M+P,\quad P\xrightarrow{\delta_p}\varnothing,\\[-2pt]
\dot m=k_m-\delta_m m,\qquad\dot p=k_pm-\delta_p p.
\end{align}
仅当 mRNA 恢复远快于蛋白质变化时，可近似令 $m\simeq k_m/\delta_m$；若研究振荡或短时间响应，该时间尺度消去可能失效。

要判断上述消去何时可靠，可以先把两阶段模型完整解出来。记两个稳态浓度为
\[
m^*=\frac{k_m}{\delta_m},\qquad p^*=\frac{k_pk_m}{\delta_p\delta_m}.
\]
分别把两个方程乘以相应的指数因子再积分，这种方法称为\emph{积分因子法}；指数因子帮助把左边合成一个乘积的导数，从而得到
\begin{align}
m(t)&=m^*+(m_0-m^*)e^{-\delta_mt},\\
p(t)&=p^*+(p_0-p^*)e^{-\delta_pt}
 +k_p(m_0-m^*)\frac{e^{-\delta_mt}-e^{-\delta_pt}}{\delta_p-\delta_m}.
\end{align}
当 $\delta_p=\delta_m=\delta$ 时，最后一项取极限为 $k_p(m_0-m^*)te^{-\delta t}$。各项均随时间减小到零，所以无论从什么非负浓度出发，两个变量最终都会回到同一个稳态。若转录率变为随时间变化的 $r(t)$，则
$m(t)=m_0e^{-\delta_mt}+\int_0^t e^{-\delta_m(t-s)}r(s)\dd s$。
如果转录率 $r(t)$ 为 $C^2$（有连续的一阶和二阶导数，即变化速度及其变化也连续），并且输入变化的时间尺度远长于 mRNA 的恢复时间，那么起始时的快速调整结束后，即 $t\gg1/\delta_m$，就可以把积分近似写成
$m(t)\simeq r(t)/\delta_m-\dot r(t)/\delta_m^2$。故仅有 $\delta_m\gg\delta_p$ 还不够，输入与其他调控变量也必须在 mRNA 时间尺度上缓慢变化。

这些速率也决定了自然的时间和浓度单位。若物理浓度 $P$ 满足 $\dot P=A+B P^n/(K_P^n+P^n)-\gamma P$，取 $x=P/K_P$、$\tau=\gamma t$，则
$\dd x/\dd\tau=a+b x^n/(1+x^n)-x$，其中 $a=A/(\gamma K_P)$、$b=B/(\gamma K_P)$。后文的阈值应在对应单位下解读。

\subsection{调控因子怎样改变产生率}
前面的产生率是常数。若调控因子与启动子结合，不同结合状态可以有不同转录率，产生率就成为调控因子浓度的函数。考虑 $D+X\rightleftharpoons DX$，令总启动子浓度 $D_T=[D]+[DX]$ 守恒。固定游离调控因子浓度 $x$ 时，平衡条件为
\[
k_{\rm on}[D]x=k_{\rm off}[DX],\qquad K_d=\frac{k_{\rm off}}{k_{\rm on}}=\frac{[D]x}{[DX]}.
\]
把守恒关系代入，得到结合概率 $p=[DX]/D_T=x/(K_d+x)$。若未结合与已结合的转录速率分别为 $r_0,r_1$，对两个状态加权就有
\begin{equation}F(x)=r_0(1-p)+r_1p=r_0+(r_1-r_0)\frac{x}{K_d+x}.\label{eq:occupancy}\end{equation}
由这个式子直接看出，$r_1>r_0$ 是激活，$r_1<r_0$ 是抑制；结合本身的存在并不能决定调控符号。这里的 $x$ 是游离浓度。若结合显著耗竭调控因子，还需要同时解其守恒关系，不能把总浓度直接代入。

这种有效产生率把启动子状态隐藏在占据概率中，因此还须检查启动子能否足够快地跟上浓度变化。固定 $x$ 时，结合概率满足
\[
\dot p=k_{\rm on}x(1-p)-k_{\rm off}p,
\]
启动子接近结合平衡的速率为 $k_{\rm on}x+k_{\rm off}$。如果结合与解离远快于表达量的变化，就可以近似认为：浓度每次改变后，启动子几乎立即重新达到相应的结合比例。这就是后文所说的“准平衡”近似。Hill 函数 $x^n/(K^n+x^n)$ 则进一步把复杂结合压缩为有效的陡峭响应；$n$ 不应直接解读为实际结合分子数。多个操纵位点、协同结合及 DNA loop 等机制都可能改变这种压缩的具体形式。

\subsection{进入回路分析之前：浓度会不会越界？}
产生率的非负性和一阶降解还有一个共同作用：确保模型始终描述有物理意义的浓度。当浓度降到 $x_i=0$ 时，有 $\dot x_i=F_i\geq0$，变化方向只能向上或暂时不变，所以浓度不会继续降成负数。在上述唯一轨迹的条件下，从非负浓度出发就始终保持非负；数学上称为非负浓度区域的\emph{正向不变性}。“正向”指随着时间向前演化，“不变”指轨迹不会离开这个区域。把方程两侧乘以 $e^{\gamma_i t}$ 再积分，还可以写出：
\[
x_i(t)=e^{-\gamma_it}x_i(0)+\int_0^t e^{-\gamma_i(t-s)}F_i(\bm x(s))\dd s.
\]
如果产生率进一步有上界 $F_i\leq M_i$，就得到
\begin{equation}
0\leq x_i(t)\leq x_i(0)e^{-\gamma_it}+\frac{M_i}{\gamma_i}(1-e^{-\gamma_it}).\label{eq:bound}
\end{equation}
这个上界说明，产生率最高为 $M_i$ 时，一阶降解会限制浓度的增长；它不会在有限时间内升到无限大，方程也能一直计算下去。后文的 Hill 产生率都有这样的上限。接下来可以据此分析系统最终停在哪里，以及不同初始状态会不会导致不同结果。

\section{一个变量：反馈怎样带来恢复或记忆？}
现在让产生率依赖蛋白质自身。单个动态变量只能沿浓度轴运动，但产生曲线的形状已经能决定系统是回到唯一表达水平，还是在不同长期状态之间作出选择。

\begin{figure}[H]\centering
\begin{tikzpicture}
\node[gene] (none) at (0,0) {$X$};\node at (0,-.8) {无反馈};
\node[gene] (pos) at (4,0) {$X$};\node at (4,-.8) {自激活};
\node[gene] (neg) at (8,0) {$X$};\node at (8,-.8) {自抑制};
\draw[act] (pos) to[loop above,min distance=12mm] (pos);
\draw[rep] (neg) to[loop above,min distance=12mm] (neg);
\end{tikzpicture}
\caption{单节点的三种符号拓扑。所有模型仍包含产生与降解；自边表示自身浓度对产生率的作用。}
\end{figure}
\subsection{自抑制加强恢复，为什么不会产生多个稳态？}
无反馈模型的产生率恒为 $\alpha$。自抑制把它改成随 $x$ 非增的 $F(x)$，例如 $F(x)=a/(1+(x/K)^n)$。浓度上升时，产生减少而降解增加，两种作用都将浓度拉回。

把这种直觉写成方程，令 $f(x)=F(x)-\gamma x$，假设产生率 $F(x)$ 非负、有界、非增且为 $C^1$。这里“有界”就是有一个有限上限，“非增”表示浓度增加时产生率不会上升，$C^1$ 则保证可以连续地计算斜率。$f$ 就严格递减，且 $f(0)\geq0$、$f(x)\to-\infty$。因此产生曲线与降解直线恰好相交一次，稳态 $x^*$ 唯一，允许它位于边界 $x^*=0$。

不仅稳态唯一，所有非负初值也都会趋于它。令浓度偏差 $e=x-x^*$，则
\[
\frac12\frac{\dd}{\dd t}e^2=e\{F(x)-F(x^*)\}-\gamma e^2\leq-\gamma e^2,
\]
因为 $F$ 非增，浓度偏差与产生率偏差的乘积不可能为正。积分可得
\[
|x(t)-x^*|\leq e^{-\gamma t}|x(0)-x^*|.
\]
式~\eqref{eq:bound}说明浓度不会无限增大，所以这个恢复速度的估计可以一直使用。若 $F'(x^*)<0$，小扰动满足 $\dot e\simeq(F'(x^*)-\gamma)e$，小偏差减小到原来约 $37\%$ 所需的时间为 $1/(\gamma-F'(x^*))$，短于无反馈的 $1/\gamma$。这把“自抑制促进恢复”落实为具体速率；它描述的是稳态附近的小偏差；输入变化很大时，还要直接计算整条时间轨迹。

\subsection{自激活让产生与降解出现三次交点}
自激活的方向相反：浓度上升也会提高产生率，偏差可能被放大。于是一个自然的问题是，这种放大能否使系统保留不同表达状态？取
\begin{equation}\dot x=a+b\frac{x^n}{K^n+x^n}-\gamma x.\label{eq:auto}\end{equation}
稳态是产生曲线与线性降解直线的交点。局部稳定条件为
\begin{equation}F'(x^*)-\gamma<0.\end{equation}
图~\ref{fig:auto} 的参数 $a=0.05,b=3.5,K=\gamma=1,n=4$ 产生三个稳态，外侧稳定，中间不稳定。关键不是正反馈必然双稳态，而是正反馈与足够的非线性产生多个交点。改变参数时，两个交点可能靠近并合在一起，再一起消失。合并的临界位置称为“折叠点”：此时产生曲线与降解直线不仅相交，斜率也相同，即 $F(x^*)=\gamma x^*$ 和 $F'(x^*)=\gamma$。
\begin{figure}[H]\centering\includegraphics[width=.74\linewidth]{figures/self_activation.pdf}\caption{自激活模型的产生率与降解率，实心点为稳定的稳态，空心点为不稳定的稳态。}\label{fig:auto}\end{figure}

先判断交点两侧的运动方向，再问交点何时会成对出现。对只有一个浓度变量的 $C^1$ 方程，若 $f'(x^*)<0$，则足够靠近稳态时 $f$ 在左侧为正、右侧为负，轨道向稳态运动；若 $f'(x^*)>0$ 则方向相反。如果交点处的斜率为零，就需要再看它左右两侧的实际变化方向，不能只凭这个斜率判断。对式~\eqref{eq:auto}，
\[
F'(x)=\frac{bnK^nx^{n-1}}{(K^n+x^n)^2}.
\]
当 $n>1$ 时该导数的唯一最大值出现在 $x_m=K((n-1)/(n+1))^{1/n}$：令 $h=x/K$，对 $h^{n-1}/(1+h^n)^2$ 求导，驻点条件为 $(n-1)-(n+1)h^n=0$。因此
\[
\max F'=\frac bK\frac{(n+1)^2}{4n}
\left(\frac{n-1}{n+1}\right)^{(n-1)/n}.
\]
若该值小于 $\gamma$，$f$ 严格递减，稳态唯一；超过 $\gamma$ 只是出现多交点的必要条件，仍需选择合适的 $a$。
令 $q=(x^*/K)^n$，联立折叠条件可参数化为
\[
b_{\rm fold}=\frac{\gamma K}{n}q^{(1-n)/n}(1+q)^2,\qquad
a_{\rm fold}=\frac{\gamma K}{n}q^{1/n}(n-1-q).
\]
这是由 $F'=\gamma$ 先解出 $b$，再代入 $a=\gamma x-bq/(1+q)$ 得到。对 $a>0$ 需 $0<q<n-1$。通常的折叠分岔还要求 $f''\ne0$，称为一个\emph{非退化条件}。“分岔”指参数变化使稳态或轨迹的类型发生改变；“非退化”在这里指相切处仍有明确的弯曲，两个交点的合并不是更特殊的平坦情形。若连弯曲程度也为零，交点怎样变化就需要进一步分析。

把条件落实到图中的参数，也能解析地确认这三个交点。取 $f(x)=0.05+3.5x^4/(1+x^4)-x$，有 $f(0)>0$、$f(0.1)<0$、$f(1)>0$、$f(4)<0$，所以三个互不相交的区间 $(0,0.1)$、$(0.1,1)$、$(1,4)$ 各含至少一个根。上面的计算说明 $f'$ 最多有两个零点，也就是说曲线最多改变两次升降方向。若 $f$ 有四个交点，相邻交点之间就至少需要三个斜率为零的位置，超过了这个上限。这用到了\emph{Rolle 定理}：一条可求导的曲线在两端取相同值时，中间至少有一处斜率为零。因此三个交点就是全部结果。两个极值分别严格低于、严格高于零，因此三个根均不在极值点处，斜率依次为负、正、负。中间的不稳定交点构成一道分界：初始浓度在它以下时趋于低表达态，在它以上时趋于高表达态。短暂扰动跨过此阈值可造成持久状态切换。

\subsection{有记忆，为什么仍不能自行振荡？}
两个稳定表达态已经足以构成记忆，却还不足以构成时钟。对\emph{一维自治方程} $\dot x=f(x)$，“一维”表示只有一个动态变量，“自治”表示变化率由当前状态决定、不显式依赖时间。如果外部条件固定、也没有延迟，那么同一个浓度 $x$ 总对应同一个变化方向。浓度在一个区间里一直上升，就不能在没有其他变化的情况下突然掉头下降。

可以把这一点完整地论证出来。假设有一条重复往返的轨迹，周期为 $T$。它在一个周期内一定有最高浓度 $x_m$。最高点处的变化率为零，所以 $f(x_m)=0$，也就是说这里已经是稳态。从这个浓度出发，始终停在 $x_m$ 就满足原方程；同一初始浓度又只能产生一条轨迹，因此原来的轨迹也只能停在那里，与“重复往返”矛盾。

所以，一个变量可以提供不同表达状态之间的选择，但在这些假设下不能自行维持振荡。要寻找持续振荡，需要增加动态环节，或者另行考虑延迟、随时间变化的输入等因素。

\section{两个变量：从信号传递到反馈开关}
增加第二个变量后，轨道从浓度轴进入平面。新节点可以只是接收上游信号，也可以将作用传回上游。先比较这两种连接，才能看出闭环符号如何影响扰动，以及二维模型是否已经足以振荡。
\subsection{只传递信号：下游跟随上游}
\begin{equation}\dot x=\alpha_x-\gamma_xx,\qquad\dot y=G(x)-\gamma_yy.\end{equation}
$x$ 的稳态先独立确定，再决定 $y$。纯前向连接不能让 $y$ 反过来影响 $x$；它改变输出幅度、响应时间和滤波性质，而不是构成闭合反馈回路。

这一连接的长期状态可以按上游、下游的顺序确定。若 $G$ 连续非负，则 $x^*=\alpha_x/\gamma_x$、$y^*=G(x^*)/\gamma_y$ 唯一。上游的小偏差按速率 $\gamma_x$ 衰减，下游自身的偏差按速率 $\gamma_y$ 衰减。要进一步确认下游最终确实跟到对应的稳态，可以直接看它的解：
\[
y(t)=e^{-\gamma_yt}y_0+\int_0^t e^{-\gamma_y(t-s)}G(x(s))\dd s.
\]
当 $x(s)\to x^*$ 时，$G(x(s))$ 也趋于 $G(x^*)$。积分中较早的输入影响会随时间衰减，较晚的输入又已接近稳态值，因此 $y(t)\to y^*$。如果上游在稳态附近有小幅、反复的变化，设重复一次所需的时间为 $T$，角频率为 $\omega=2\pi/T$，那么上游变化 $\delta x$ 与下游变化 $\delta y$ 的幅度和时间滞后关系为
$\delta y(i\omega)/\delta x(i\omega)=G'(x^*)/(\gamma_y+i\omega)$：高频幅值被抑制，相位发生滞后。这个关系把上游浓度变化作为输入；若改变的是上游产生率，则还要考虑上游浓度本身需要多久才能跟上。

\subsection{闭合回路：反馈符号怎样进入稳定性？}
最简单的无自反馈模型是
\begin{equation}\dot x=f(y)-\gamma_xx,\qquad\dot y=g(x)-\gamma_yy.\label{eq:two}\end{equation}
两条边可以是 $++$（互激活）、$--$（互抑制）、$+-$ 或 $-+$（一激活一抑制）。要知道轻微偏离稳态后会怎样，可以对方程作\emph{线性化}，也就是用稳态附近的斜率近似曲线，只保留小偏差的一阶影响，把它们整理成下面的矩阵 $J$，通常称为雅可比矩阵。对角项表示每个变量自身的恢复作用，另外两项表示一个变量怎样影响另一个：
\begin{equation}J=\begin{pmatrix}-\gamma_x&f'(y^*)\\g'(x^*)&-\gamma_y\end{pmatrix},\quad \operatorname{tr}J=-\gamma_x-\gamma_y<0,\quad\det J=\gamma_x\gamma_y-f'g'.\end{equation}
正反馈时 $f'g'>0$，两个节点会沿回路放大某些偏差。当 $f'g'>\gamma_x\gamma_y$，这股放大作用超过降解的恢复作用，至少有一种偏差会越变越大；另一个方向的偏差仍会恢复，这样的稳态叫“鞍点”。是否还存在另外两个稳定状态，需要继续寻找所有产生与降解平衡的位置。负反馈时 $f'g'<0$，这个模型的稳态附近偏差会衰减，有时还会边衰减边上下摆动。

这里判断的是偏差随时间放大还是缩小。偏差可以分成两种基本变化方式，每一种都带有 $e^{\lambda t}$ 这样的时间因子，$\lambda$ 称为特征值。$\lambda<0$ 表示这一种偏差衰减，$\lambda>0$ 表示它增长；若 $\lambda$ 是复数，偏差还会在增长或衰减的同时摆动，增长还是衰减由它的实部决定。令 $q=f'(y^*)g'(x^*)$，把两种变化方式的特征值写出来，得到
\[
\lambda_{\pm}=-\frac{\gamma_x+\gamma_y}{2}
\pm\sqrt{\frac{(\gamma_x-\gamma_y)^2}{4}+q}.
\]
这两个数满足 $\lambda^2+(\gamma_x+\gamma_y)\lambda+\gamma_x\gamma_y-q=0$。前面写出的 $\operatorname{tr}J$ 称为\emph{迹}，是对角项之和；$\det J$ 称为\emph{行列式}，对这里的二维矩阵就是两个对角项的乘积减去两个交叉项的乘积。它们分别对应两个特征值的和与积，因此能用来判断偏差会怎样变化。若 $\det J>0$，两个实数特征值的和为负、积为正，所以都为负；若它们是复数，共同的实部为 $-(\gamma_x+\gamma_y)/2$，也为负。两种情况都表示偏差会衰减。若 $\det J<0$，两个特征值一正一负，出现鞍点。若 $\det J=0$，至少有一个特征值为零，仅看稳态附近的斜率还不能判断这个方向会怎样。
负反馈的衰减振荡条件进一步是
$q<-(\gamma_x-\gamma_y)^2/4$，所以负反馈符号本身甚至不保证\emph{衰减}振荡；正反馈 $q>0$ 则两种变化方式都不会带有这种摆动。
\begin{figure}[H]\centering
\begin{tikzpicture}
\foreach \i/\title in {0/互激活,1/互抑制,2/负反馈}{
\node at (4.1*\i+1,1.2) {\title};
\node[gene] (x\i) at (4.1*\i,0) {$X$};
\node[gene] (y\i) at (4.1*\i+2,0) {$Y$};}
\draw[act] (x0) to[bend left=25] (y0);\draw[act] (y0) to[bend left=25] (x0);
\draw[rep] (x1) to[bend left=25] (y1);\draw[rep] (y1) to[bend left=25] (x1);
\draw[act] (x2) to[bend left=25] (y2);\draw[rep] (y2) to[bend left=25] (x2);
\end{tikzpicture}
\caption{两个交叉边形成的反馈。互激活与互抑制均为正回路；一激活一抑制为负回路。交换节点得到另一种带标签的负回路。}
\end{figure}

在负反馈下，局部稳定还可以与稳态的唯一性联系起来。假设 $f,g$ 非负、有界且为 $C^1$，并且一个随输入增加而严格上升，另一个严格下降。在稳态方程中消去 $y$，得到
\[
x=T(x):=\frac1{\gamma_x}f\!\left(\frac{g(x)}{\gamma_y}\right).
\]
经过一正一负两次调控，$T$ 不会随 $x$ 增加而上升，而且有上限。所以 $T(x)-x$ 一直下降，在零处非负，$x$ 足够大时为负，恰好经过零一次。随后 $y=g(x)/\gamma_y$ 也唯一。又有 $q\leq0$，由上面的计算，两种偏差都只能衰减，因此这个唯一稳态在小偏差后能逐渐恢复。这与单节点自抑制相呼应：闭环总体为负时，偏差受到恢复作用；两个变量则额外允许恢复过程带有衰减振荡。

\subsection{二维相图为什么仍没有持续振荡？}
边衰减边摆动，仍然是逐渐回到稳态；持续振荡则要求轨迹每隔一段时间回到原处，并重复同样的变化。对式~\eqref{eq:two}，可以排除后一种情况，原因是这个模型始终使浓度平面中的面积收缩。这是\emph{Bendixson 判据}在此模型中的应用：若一个没有空洞的平面区域内，散度始终严格为同一符号，就不存在非恒定周期轨道。下面解释为什么面积收缩会产生这个限制。

把两个变化率写成 $\bm v=(f(y)-\gamma_xx,g(x)-\gamma_yy)$，则沿各自浓度方向计算斜率并相加，有
\[
\frac{\partial\dot x}{\partial x}+\frac{\partial\dot y}{\partial y}=-\gamma_x-\gamma_y<0.
\]
这个和称为“散度”。它为负，表示一小片初始状态经过同样一段时间的演化后，占据的面积会缩小。

假如存在一条不断重复的轨迹，它在浓度平面上就围成闭合曲线 $C$，内部区域记为 $\Omega$。轨迹不会自相交，因为同一组浓度只能对应一种后续变化。还假设这条曲线及其内部都处于方程可以正常使用的范围内。数学上常通过\emph{单连通区域}保证这一点：直观地说，区域中没有空洞，区域内的闭合曲线所围住的内部也属于这个区域。沿曲线，变化方向与曲线相切，没有穿过边界向外的净流动；但内部处处收缩又要求净流动小于零。平面散度定理把“沿边界的净流动”与“内部收缩率的总和”联系起来，因此这两种说法可以写成同一个面积关系：
\[
0=\oint_C\bm v\cdot\bm n\,\dd\ell
=\iint_\Omega\nabla\cdot\bm v\,\dd x\dd y
=-(\gamma_x+\gamma_y)\operatorname{Area}(\Omega)<0,
\]
相互矛盾。左边沿边界计算净向外流动，右边把内部的收缩率加起来；$\bm n$ 表示指向边界外侧的方向。因此该模型没有持续重复的非恒定轨迹。

这里用到了“没有自调控”和“一阶降解”。换成其他二维负反馈模型，收缩率未必始终为负，结论也可能改变。另外，没有持续振荡并不等于所有初始状态都回到同一个稳态：正反馈仍可能形成多个长期状态。

\subsection{互抑制怎样把表达状态分成两种选择？}
负反馈倾向于恢复，正反馈则可能放大两个节点之间的差异。互抑制尤其直观：$X$ 增多会压低 $Y$，较低的 $Y$ 又解除对 $X$ 的抑制，最初的表达优势于是被增强。用下面的无量纲模型分析这一过程：
\begin{equation}\dot x=\frac{a}{1+y^n}-x,\qquad\dot y=\frac{a}{1+x^n}-y.\label{eq:toggle}\end{equation}
先把两条变化率分别设为零，得到 $x=a/(1+y^n)$ 与 $y=a/(1+x^n)$。这样的曲线称为“零流线”：在一条线上 $x$ 暂时不变，在另一条线上 $y$ 暂时不变；两条线的交点处两个变量都不变，就是稳态。图~\ref{fig:toggle} 展示 $a=3,n=3$ 的示例：两个不对称表达态和中间的对称稳态。再检查这些交点附近的小偏差会恢复还是被放大，就能判断它们是否稳定。该模型展示状态选择的最小机制，但生物学中的细胞命运并不等同于某一对基因浓度。

先从两个基因表达相同的状态开始，问微小的不对称是否会增长。令 $x=y=s$，则 $a=s(1+s^n)$；右侧在 $s\geq0$ 严格递增，故对称稳态唯一。记交叉斜率的绝对值为
$h=ns^n/(1+s^n)$，雅可比为 $\begin{psmallmatrix}-1&-h\\-h&-1\end{psmallmatrix}$。
如果两个浓度同时增加同样的一小点，偏差按 $e^{(-1-h)t}$ 衰减；如果一个增加、另一个减少同样的一小点，差异则按 $e^{(-1+h)t}$ 变化。这分别对应特征值 $-1-h$ 与 $-1+h$。当 $n>1$ 时，后者过零的阈值为
\[
s_c=(n-1)^{-1/n},\qquad a_c=\frac n{n-1}(n-1)^{-1/n}.
\]
当 $n\leq1$ 时有限 $s>0$ 下 $h<1$，对称态不会经这一线性机制失稳。阈值只判断对称态；不对称态的存在和稳定性仍须证明。

对于图中的 $a=3,n=3$，可以继续把所有稳态解出来，从而确认系统有一个对称鞍点和两个稳定的不对称表达态。
稳态满足 $x+xy^3=3$、$y+yx^3=3$，故 $x,y>0$。相减得
$(x-y)[1-xy(x+y)]=0$。对称解满足 $s+s^4=3$，严格单调性保证唯一；又 $s>1$，故 $h=3s^3/(1+s^3)>1$，是鞍点。
不对称时令 $S=x+y$、$P=xy$，有 $PS=1$。相加得
$S+P(S^2-2P)=6$，即 $S-S^{-2}=3$。左侧在 $S>0$ 严格递增，故恰有一个解，且 $S>3$。二次方程 $t^2-St+1/S=0$ 的判别式 $S^2-4/S>0$，两根均正；交换两根得恰好两个不对称解。此和、差条件亦反推出原两条稳态方程，未引入额外解。
在这些点，两个交叉斜率均等于 $-3x^2y^2/a=-P^2$，特征值为 $-1\pm P^2$。因为 $P=1/S<1/3$，两根均负。三种分支已穷尽所有稳态。小的不对称在鞍点附近被放大，而两个不对称态附近的偏差又被压回去；这两步共同构成了开关的状态选择与保持。

\begin{figure}[H]\centering\includegraphics[width=.67\linewidth]{figures/toggle_nullclines.pdf}\caption{互抑开关的零流线与相流（$a=3,n=3$）。实心点稳定，空心点为鞍点；绿色轨迹分别趋于两种不对称表达态。}\label{fig:toggle}\end{figure}

\subsection{把自反馈加回来：怎样完整整理 81 种情况？}
前面的模型只保留两条交叉边。允许两个节点分别自激活、自抑制或没有自边后，产生率变为 $F_x(x,y),F_y(x,y)$，一般形式是
\begin{equation}
\dot x=F_x(x,y)-\gamma_xx,\qquad
\dot y=F_y(x,y)-\gamma_yy.\label{eq:two_general}
\end{equation}
以下分类假设产生率非负、有界且为 $C^1$，边的符号在所讨论的浓度范围内固定；缺失边表示相应偏导数恒为零。输入固定，没有延迟，也没有其他隐藏动态变量。这样的约定让“一个拓扑”具有可比较的含义；若边会随浓度改变符号，就不属于同一个固定符号类别。

四条边 $(X\to X,X\to Y,Y\to X,Y\to Y)$ 各有 $0,+,-$ 三种选择，所以带标签的情况恰为 $3^4=81$。交换 $X,Y$ 后不变的图需要两条自边相同、两条交叉边相同，共 $3^2=9$ 个；其余图两两配对，故有 $9+(81-9)/2=45$ 个无标签类别。本文先保留标签，方便区分输入、输出及两个节点的不同时间尺度。

逐一背 81 张图不利于理解。更有效的办法是先看两个节点是否互相连通，再看有没有自激活：没有交叉边有 $9$ 种；恰有一条交叉边有 $2\times2\times9=36$ 种；两条交叉边都存在时，同号与异号各有 $2\times9=18$ 种。每种连接下，自边的 $9$ 种组合中有 $4$ 种完全没有正自边、$5$ 种至少有一条正自边。因此所有情况恰好分成下表的八组。

\begin{center}\small
\begin{tabular}{c r >{\raggedright\arraybackslash}p{4.6cm} >{\raggedright\arraybackslash}p{7.6cm}}
\toprule 组 & 数目 & 结构 & 先分析什么\\\midrule
R1 & 4 & 无交叉边，无自激活 & 两个独立的恢复过程；每个节点有唯一稳态。\\
R2 & 5 & 无交叉边，有自激活 & 独立记忆单元；自激活足够强时可有多个表达组合。\\
R3 & 16 & 单向连接，无自激活 & 有限增益的信号传递；唯一稳态，不能自行持续振荡。\\
R4 & 20 & 单向连接，有自激活 & 上游状态传给下游，或下游自己存储状态；先后求解两个节点。\\
R5 & 8 & 同号交叉边，无自激活 & 互激活或互抑制开关；多稳态取决于闭环增益。\\
R6 & 10 & 同号交叉边，有自激活 & 多个正反馈共同作用；考察状态数、阈值与迟滞，不能经 Hopf 失稳。\\
R7 & 8 & 异号交叉边，无自激活 & 唯一局部渐近稳定的稳态，可有衰减摆动，无周期轨道。\\
R8 & 10 & 异号交叉边，有自激活 & 正放大与负恢复竞争；可能出现 Hopf 分岔，也可能只是恢复或多稳态。\\
\bottomrule
\end{tabular}\end{center}

这些是\emph{结构分组与解析筛选}，不是把每张图标成固定功能。R1 的恢复看起来简单，却可以提供稳定的基线；R3 可以滤去快速噪声；R2、R4 的记忆也不必依赖跨节点闭环。哪些情况“有趣”，取决于想解释的是恢复、状态选择、时间响应还是振荡。图~\ref{fig:81map}把 81 个组合完整放进同一张索引图；附录还给出逐张回路图与编号，方便查到任何一种情况。
\begin{figure}[H]\centering
\includegraphics[width=.76\linewidth]{figures/topology_classes.pdf}
\caption{完整的 81 种两节点符号拓扑。行是 $(X\to X,Y\to Y)$ 的自边组合，列是 $(X\to Y,Y\to X)$ 的交叉边组合；$0$ 表示缺失。每个格子的 R 编号对应正文八组，颜色表示分组，不表示已经实现某种功能。}\label{fig:81map}
\end{figure}

\subsection{先排除简单情况，再识别记忆的来源}
R1 与 R2 没有节点间连接，整个系统是两个一维系统并列。R1 的自边只有缺失或抑制，前面的一维恢复论证分别适用。R2 则允许其中一个或两个节点具有自激活；例如两个节点都在双稳态参数区时，每个节点各有低、高两种稳定状态，组合后可以有四种稳定表达状态。这是两份独立记忆，并不需要互抑制。

R3 与 R4 虽然有传递，但作用只有一个方向。先求不受下游影响的节点，再把它的稳态代入下游方程。R3 的两条自边都不会促进自身产生，因此两个步骤各有唯一稳态；上游趋于稳态后，下游也会逐渐跟上。严格地看，下游偏差的恢复率至少为 $\gamma_y$，而上游尚未稳定造成的附加输入会趋于零，所以前面的积分估计仍保证它最终恢复。R4 的任一步都可能因自激活而出现多个解：需要分别问“上游能否选择状态”和“同一上游输入下，下游能否保留不同状态”。这两种记忆机制值得区分。

这四组总共有 $45$ 种，都没有跨节点闭环。它们不能在固定输入下产生非恒定周期轨道：若整个系统周期变化，上游也必须周期变化；一维自治方程迫使上游恒定，随后下游又成为一个一维自治方程，也只能恒定。因此“有两个变量”本身不是振荡机制。

对 R5 与 R6，同号交叉边构成正回路。互激活增强共同表达的变化，互抑制则增强两者之间的差异；它们与自激活一样，都能为状态选择提供放大作用，但还需要足够陡的产生率曲线。一个有用的结构统计是：含至少一条正自边的图有 $81-2^2\times3^2=45$ 种；没有正自边却有正交叉回路的图有 $2\times4=8$ 种，所以含正回路的共有 $53$ 种。

其余 $28$ 种就是 R1、R3、R7。前两组已知只有一个稳态，下面说明 R7 也如此。因此在本文的固定符号、有限产生率模型中，多稳态只能在这 $53$ 种含正回路的图里寻找。这个条件只是筛选入口，不代表这 $53$ 种在任意参数下都会多稳态。

R7 可以保留自抑制。给定 $y$ 时，$F_x(x,y)-\gamma_xx$ 随 $x$ 严格下降，恰有一个零点 $x=h(y)$；同理 $y=k(x)$ 唯一。由定态方程求导，
\[
h'(y)=\frac{\partial_yF_x}{\gamma_x-\partial_xF_x},\qquad
k'(x)=\frac{\partial_xF_y}{\gamma_y-\partial_yF_y}.
\]
分母均为正，两条交叉边异号，所以 $h(k(x))$ 非增，与直线 $x$ 只相交一次。边界处的解也可以由同样的单调性处理。因而 R7 有唯一稳态；两个对角恢复率为负、交叉作用乘积为负，保证它局部渐近稳定。散度在整个区域为负，又排除了周期轨道。可见无反馈、负反馈和弱正反馈都可能表现为恢复，但原因并不相同。

\subsection{哪些自反馈真正打开了二维振荡的入口？}
在稳态附近，式~\eqref{eq:two_general}的雅可比可以写成
\[
J=\begin{pmatrix}a&b\\c&d\end{pmatrix},\quad
 a=\partial_xF_x-\gamma_x,\quad d=\partial_yF_y-\gamma_y,
\quad b=\partial_yF_x,\quad c=\partial_xF_y.
\]
这里 $a,d$ 是自身产生与降解合起来的\emph{净作用}，不是自边符号本身。正自边也可能太弱，以至于 $a$ 或 $d$ 仍为负。两个特征值满足
\[
\lambda_{\pm}=\frac{a+d}{2}\pm\frac12\sqrt{(a-d)^2+4bc}.
\]
这一个公式便能检验全部八组的 Hopf 入口。R1--R4 至少一条交叉边缺失，$bc=0$，两根为实数；R5、R6 有 $bc>0$，根号内严格为正，两根也为实数，所以不会通过一对复特征值穿越虚轴产生 Hopf 分岔。R7 虽然 $bc<0$，但 $a,d<0$，迹不可能到零；没有正自边的所有 $36$ 种图还都满足全域面积收缩。

于是只有 R8 的 $10$ 种具备这种失稳的符号条件：两条交叉边异号，而且至少有一条正自边。在候选临界点，
\[
a+d=0,\qquad ad-bc>0.
\]
代入 $d=-a$，条件变成 $-bc>a^2$。直观地说，一边的自放大刚好抵消另一边的净恢复，而负闭环还必须足够强，才能使偏差以摆动方式变化。这仍只是必要的线性条件，真正的 Hopf 分岔还需参数横截和非线性非退化条件。

为了看到“同一结构可以恢复，也可以振荡”，考虑 R8 中的一张图：$X$ 自激活，$X$ 激活 $Y$，$Y$ 抑制 $X$，$Y$ 没有自边。取
\begin{align}
H(x)&=\frac{x^4}{1+x^4},\\
\dot x&=\frac{a_0+b_0H(x)}{1+y^4}-x,\qquad
\dot y=\varepsilon\{d_0H(x)-y\}.
\end{align}
自激活让 $X$ 快速上升，较慢积累的 $Y$ 随后压低 $X$；$X$ 降低后，$Y$ 的产生也减少，抑制逐渐解除。这就是正放大与负恢复交替工作的机制。固定 $a_0=0.4,b_0=4,d_0=2$，图~\ref{fig:twoosc}比较 $\varepsilon=1$ 与 $\varepsilon=0.05$：前者在小扰动后恢复，后者数值上持续振荡。改变的是 $Y$ 的时间尺度，而不是连边。脚本检查了示例定态的局部稳定性及后期轨迹，图中振荡仍是数值例子，没有替代周期轨道的严格证明。
\begin{figure}[H]\centering\includegraphics[width=.95\linewidth]{figures/two_node_oscillation.pdf}
\caption{同一 R8 拓扑的两种参数行为。左：下游较快，偏差衰减；右：下游较慢，自激活与负反馈轮流主导，数值轨迹持续摆动。仅展示后期轨迹，未标注不必要的数值根。}\label{fig:twoosc}
\end{figure}

二维部分的结论可以这样使用：先凭结构排除不适合的机制，再用反馈增益与时间尺度检验剩余候选。开关重点看交点与正回路；Hopf 振荡重点看 R8、迹与行列式。增加第三个节点，则可以在没有正自边的负反馈环中提供新的动态环节，也可以让输入沿多条路径抵达输出。

\section{从开关到可调时钟：Science（2008）中的正负反馈}
前面的 R8 例子已经说明，正反馈不只用于记忆，还能与负反馈一起产生振荡。但“有振荡”与“适合做时钟”仍是两个不同的问题：改变周期时，峰值会不会一起消失？参数有一些变化时，振荡是否还能维持？

Tsai、Choi、Ma、Pomerening、汤超与 Ferrell 的 \emph{Robust, Tunable Biological Oscillations from Interlinked Positive and Negative Feedback Loops}（Science, 2008，参考文献~8）研究了这一问题。他们在所比较的细胞周期及简化振荡模型中发现，正负反馈交织可支持较宽的频率调节范围，同时维持近似不变的振幅；随机参数比较也显示了较大的振荡参数区域。这是模型之间的性能比较，不是“任意网络加一条正边都会更好”的定理。

\subsection{评价振荡：周期、振幅与参数范围要分别看}
设相邻两个峰的时间间隔为 $T$，频率为 $f=1/T$，一周期内的最大、最小浓度为 $x_{\max},x_{\min}$，振幅在本文记为
\[
A_{\rm osc}=x_{\max}-x_{\min}.
\]
比较振荡器时，至少要分开三个问题：它能否持续振荡，频率能调多少，以及调频时振幅改变多少。一个只能在很窄的参数区间里产生微弱摆动的系统，与一个可以调节周期、又保持明确高低状态的系统，虽然都能画出周期轨迹，功能上并不相同。

也要说明调的是哪个参数。如果把\emph{所有}速率统一乘以同一系数，等于把整个时间坐标一起缩放，当然可以改变周期而保留轨迹形状。更有意义的问题是，只改变一个合适的产生率或时间尺度，其他过程不作同步调整时，是否仍能稳定调频。稳健性则是另一个量：它考察某个采样范围内，多少组参数产生符合标准的振荡；数值取决于模型、采样规则和振荡判据。

\subsection{迟滞开关与慢负反馈，怎样把“高度”和“等待时间”分开？}
论文以细胞周期模型说明这种组合：CDK1 的正反馈包括激活 Cdc25、抑制 Wee1；经 APC 作用的负反馈则帮助结束高活动阶段。正反馈使切换更明确，负反馈使系统可以重新进入下一轮，而不是一直停在高状态（参考文献~8）。

为了把这个解释接到现有方程上，继续看本文自己的两变量模型，而不把它当作论文细胞周期模型的复现：
\[
\dot x=f(x,y),\qquad\dot y=\varepsilon g(x,y),
\]
其中 $f$ 与 $g$ 就是上一节的 R8 产生--降解规则。$\varepsilon$ 较小时，$x$ 快，抑制变量 $y$ 慢。暂时固定 $y$，只让 $x$ 恢复，就能画出\emph{快变量稳态曲线} $f(x,y)=0$。如果某个 $y$ 区间内存在明显分开的低表达和高表达两个稳定分支，中间夹着不稳定分支，$x$ 会在慢变量推动下沿一个稳定分支走，到折叠点才跳到另一支。

这就产生\emph{迟滞}：从低态转向高态与从高态转向低态，发生在不同的 $y$ 阈值。$y$ 增大时，高表达态一直保持到其分支终止，再快速降到低态；$y$ 减小时，低态又保持到另一个阈值，再快速升高。慢变化与快跳跃轮流发生，称为\emph{弛豫振荡}。这里“弛豫”是这一类振荡的标准名称，不表示整个过程都以单个指数恢复。

设高分支上 $y$ 从 $y_-$ 增到 $y_+$，低分支上再从 $y_+$ 降到 $y_-$，快分支上的稳态值分别记为 $x_H(y),x_L(y)$。若快慢分离足够明显、快跳跃所占时间很短，那么由 $\dd t=\dd y/[\varepsilon g(x,y)]$ 得到周期的近似：
\[
T\simeq\frac1\varepsilon\left[
\int_{y_-}^{y_+}\frac{\dd y}{g(x_H(y),y)}+
\int_{y_+}^{y_-}\frac{\dd y}{g(x_L(y),y)}\right].
\]
高分支上 $g>0$、低分支上 $g<0$，所以两项时间都为正。这说明慢变量的速度主要控制等待多久，而分支位置与跳跃阈值主要控制浓度跨越的范围。只改变 $\varepsilon$ 时，快稳态曲线本身不变，因此在这一近似适用的区间内，有可能明显改变周期而只小幅改变振幅。

\begin{figure}[H]\centering
\includegraphics[width=\linewidth]{figures/positive_feedback_tuning.pdf}
\caption{本文 R8 模型的快慢机制演示。左：快变量稳态曲线的稳定分支、不稳定分支与数值轨迹；中：改变同一个慢时间尺度，周期不同而浓度跨度接近；右：频率与振幅相对参考参数的变化。固定 $a_0=0.4,b_0=4,d_0=2$，仅改变 $\varepsilon$。本图不是 Science 论文原模型的数据。}\label{fig:positive_tuning}
\end{figure}

这个近似有明确限度。快慢分离减弱、参数靠近分支消失的位置，或反馈改变了快稳态曲线本身时，振幅与周期仍可能一起变化。新增脚本在所示参数范围内检查了后期重复峰、正浓度与频率变化；这些数值检查说明例子怎样工作，并不替代极限环存在性的证明。

\subsection{与 Cell 适应论文连接：正反馈的功能取决于回路角色}
这一讨论把正反馈的三种角色串起来：单独形成足够强的放大时，它可以提供多个稳定状态；与较慢负反馈配合时，它可以让系统在状态之间反复切换；在适应模型的特殊速率条件下，它还可能参与累计误差的调节。

Science（2008）关注的是时钟的可调性与振荡稳健性，后面 Cell（2009）关注的是响应之后能否回到基线。相似的正、负边可能被这两种功能以不同方式使用，评价标准也不同。因此，增加第三个节点时，我们不仅问是否更容易振荡，也问它是否能提供时间筛选、比例抵消或误差缓冲等新的动态角色。

\section{三个调控环节：持续振荡与短暂适应}
双节点已经提供两条路线：正回路可以形成开关，负回路配上正自反馈则有二维振荡的入口。若没有正自边，前面的瞬时产生--一阶降解模型仍排除了周期轨道。第三个环节提供另一种时间结构：在负闭环中，它增加动态环节，让没有正自边的模型也能出现摆动失稳；在前馈中，它让同一输入沿快慢两条路径到达输出。下面分别讨论振荡、时间筛选与适应。
\begin{figure}[H]\centering
\begin{tikzpicture}
\node[gene] (a) at (0,1.4) {$X_1$};\node[gene] (b) at (3,1.4) {$X_2$};\node[gene] (c) at (1.5,-.5) {$X_3$};
\draw[rep] (a)--(b);\draw[rep] (b)--(c);\draw[rep] (c)--(a);
\node[gene] (u) at (6,1.4) {$u$};\node[gene] (y) at (9,1.4) {$y$};\node[gene] (z) at (7.5,-.5) {$z$};
\draw[act] (u)--(y);\draw[act] (u)--node[left]{慢支路}(z);\draw[rep] (z)--(y);
\end{tikzpicture}
\caption{左：三次抑制构成负反馈环。右：不一致前馈回路，直接激活与间接抑制在输出处竞争；输入 $u$ 在这里是外部信号。}
\end{figure}

\subsection{闭合负反馈：三个环节怎样使稳态失稳？}
三个相互抑制环节可写为（本节取相同的一阶降解率 $\gamma$）
\begin{equation}\dot x_1=f_1(x_3)-\gamma x_1,\quad\dot x_2=f_2(x_1)-\gamma x_2,\quad\dot x_3=f_3(x_2)-\gamma x_3.\end{equation}
反馈符号仍然是负，但稳定性已经不能照搬双节点的结论。要找出新的可能，先用稳态附近的斜率近似产生率曲线，计算三种基本偏差各自会增长还是衰减。记 $a_i=f_i'$ 为三个交叉斜率。三次抑制使 $a_1a_2a_3<0$，写成 $a_1a_2a_3=-k^3$、$k>0$，则
\[
\lambda I-J=\begin{pmatrix}\lambda+\gamma&0&-a_1\\-a_2&\lambda+\gamma&0\\0&-a_3&\lambda+\gamma\end{pmatrix},
\quad \det(\lambda I-J)=(\lambda+\gamma)^3-a_1a_2a_3.
\]
当乘积为 $-k^3$ 时，取 $-1$ 的三个立方根 $-1,\frac12\pm i\frac{\sqrt3}{2}$ 得到
\[
\lambda_1=-\gamma-k,\qquad\lambda_{2,3}=-\gamma+\frac k2\pm i\frac{\sqrt3k}{2}.
\]
$k<2\gamma$ 时，三种偏差的增长率都为负，会逐渐减小；$k>2\gamma$ 时，其中两种带摆动的偏差开始增长。门槛 $k=2\gamma$ 处，摆动的角频率为 $\sqrt3\gamma$，对应周期 $2\pi/(\sqrt3\gamma)$。
对对称 Hill 环 $f(s)=\alpha/(1+s^n)$，唯一对称稳态满足 $\alpha=\gamma s(1+s^n)$，因此
\[
\frac k\gamma=\frac{ns^n}{1+s^n}.
\]
这里也不会漏掉另外的不对称稳态：当三个调控函数相同且都严格下降时，沿环走过三次调控，浓度增加最终对应的是反馈输入下降。这条反馈关系与 $x$ 本身只会相交一次，而三个浓度相等的解已经满足关系，因此它就是整个环的唯一稳态。
当 $n\leq2$ 时不可能达到有限参数的阈值；当 $n>2$ 时
\[
s_c=\left(\frac2{n-2}\right)^{1/n},\qquad
\alpha_c=\gamma\frac n{n-2}\left(\frac2{n-2}\right)^{1/n}.
\]
在 $n=3,\gamma=1$ 下 $\alpha_c=3\sqrt[3]2$，演示的 $\alpha=20$ 位于稳态失稳一侧。
随 $\alpha$ 增加，$\dd s/\dd\alpha=1/[\gamma(1+(n+1)s^n)]>0$，且 $\dd k/\dd s>0$；所以随参数增大，这两种摆动方式确实会从衰减转为增长。这对复特征值的实部从负变正，数学上称为\emph{穿过虚轴}；“虚轴”就是复数平面中实部为零的位置。因为这里增长率随参数变化确实不为零，这次穿越是\emph{横截}的，也就是实际穿过去，而不是碰到边界后折返。这是 Hopf 分岔的候选点：稳态附近原本衰减的摆动可能变成周期振荡。

要进一步用标准 Hopf 分岔定理判断附近的周期解，还需模型足够光滑，并检验非线性的非退化条件。常用的量是\emph{第一 Lyapunov 系数}，它描述微小振幅开始增长后，非线性会限制这种增长还是进一步增强它。标准的非退化条件要求这个系数不为零；它的符号还帮助判断新出现的周期解是否稳定。这个系数在本文尚未计算，所以不能据此宣称已经证明稳定的周期振荡。以上计算只确定了小偏差不再衰减的门槛。图~\ref{fig:rep}采用 $f(s)=20/(1+s^3)$ 展示给定初值下的持续振荡；周期轨道的存在与稳定性属于进一步的非线性问题，不能由这幅图替代证明。


\begin{figure}[H]\centering\includegraphics[width=.78\linewidth]{figures/repressilator.pdf}\caption{三节点负反馈环的数值轨迹，生成脚本见 \texttt{code/make\_figures.py}。}\label{fig:rep}\end{figure}

\subsection{快慢前馈：输出怎样先响应再回落？}
输入 $u$ 可以直接激活输出 $y$，又通过慢变量 $z$ 抑制 $y$：
\begin{equation}\tau_y\dot y=H(u,z)-y,\qquad\tau_z\dot z=G(u)-z.\end{equation}
在 $H_u>0,H_z<0,G'>0$ 的调控范围内，若 $z$ 比 $y$ 慢，上升阶跃可使 $y$ 先增加，随后 $z$ 增加抑制 $y$，形成脉冲。\emph{脉冲}与\emph{精确适应}不同：精确适应要求稳态输出 $y^*(u)$ 与持续输入的幅度无关，这对 $H$、$G$ 有额外条件；拓扑本身不保证这种性质。后面的 Cell 论文讨论会把这一支路抵消机制与带缓冲节点的负反馈作具体比较。

要把“回落”与“回到同一基线”区分开，先算恒定输入下的长期输出。设 $\tau_y,\tau_z>0$，$H,G$ 在所考察输入及状态范围内连续，固定 $u$ 时有
\[
z(t)=G(u)+(z_0-G(u))e^{-t/\tau_z},
\]
所以 $z\to G(u)$。此时输入输出支路 $H(u,z(t))\to H(u,G(u))$，按前面单向调控的方法把 $y$ 的方程积分，就得到
\[
y^*(u)=H(u,G(u)).
\]
精确适应要求输入改变后仍回到相同的输出，因此恰好要求这个复合函数在输入区间内为常数。如果 $H,G$ 为 $C^1$，且所讨论的输入覆盖一个完整区间 $I$，那么可以用\emph{链式法则}沿两条支路计算输入变化的影响。链式法则就是：间接支路的总变化率等于沿路各步变化率的乘积。由此得到
\[
H_u(u,G(u))+H_z(u,G(u))G'(u)=0\qquad(u\in I).
\]
这个式子揭示了适应所需的抵消：直接输入对输出的长期增益，与输入通过 $z$ 产生的长期增益大小相等、符号相反。快慢差异帮助形成瞬态峰，但只有这一稳态抵消才使基线与输入幅度无关。

可以选一个简单函数把这两个层面同时展示出来。取 $u>0,z>0$、$G(u)=u$、$H(u,z)=u/z$。每个恒定输入的稳态为 $(y^*,z^*)=(1,u)$，故确有精确适应。若从输入 $u_0$ 的稳态出发，在 $t=0$ 切换至 $u_1>u_0$，则
\[
z(t)=u_1+(u_0-u_1)e^{-t/\tau_z},\qquad
\dot y(0^+)=\frac{u_1/u_0-1}{\tau_y}>0.
\]
$y$ 起初升高，最终回到 $1$，所以至少存在一个高于基线的峰；$\tau_z\gg\tau_y$ 使脉冲更易分辨。$u/z$ 是理想化的比值调控，只在 $u,z>0$ 的工作域内使用。若换成 $H=u/(1+z)$，保持同一拓扑与 $G$，则 $y^*=u/(1+u)$ 依赖输入，精确适应消失。这给出了拓扑不充分的明确反例。
\begin{figure}[H]\centering\includegraphics[width=.85\linewidth]{figures/feedforward_adaptation.pdf}
\caption{相同的不一致前馈结构可以精确适应，也可以只产生瞬态峰而改变最终基线。输入从 $1$ 跳至 $2$，$\tau_y=0.25,\tau_z=2$；右图比较两种调控函数的稳态输入输出曲线。}\label{fig:adapt}
\end{figure}

\subsection{一致前馈也值得研究：它可以筛选信号持续时间}
前馈回路有两条从输入到输出的路径。若两条路径的累计符号相同，称为\emph{一致前馈}；相反则为\emph{不一致前馈}。前面的例子利用两条路径相反的长期作用形成适应，而一致前馈的一种用途是区分短脉冲与持续输入。

看一个全激活的三节点回路 $A\to B\to C$、$A\to C$。为单独看时间逻辑，暂时把输入节点 $A$ 的活动当成 $0$ 或 $1$，令 $\tau_B\dot B=A-B$，再以阈值 $0<\theta<1$ 判断 $B$ 是否足够高。如果输出产生采用 AND 逻辑，要求 $A$ 与 $B$ 同时高：上升输入到来后，$B(t)=1-e^{-t/\tau_B}$，只有等到
\[
t_{\rm on}=-\tau_B\ln(1-\theta)
\]
才开始产生输出；输入一撤去，AND 条件立即失效。所以比这个时间还短的输入可能被过滤掉。如果改成 OR 逻辑，输入到来时直接支路立即打开，撤去后则要等 $B$ 从高值降到阈值，关闭的等待时间为 $t_{\rm off}=\tau_B\ln(1/\theta)$。

这是理想阈值模型，真实的连续产生率会把“立即”与“等待”平滑成不同快慢的响应，输出自身的恢复时间也会再加一层滞后。重点是：同样的三条激活边，配上不同启动子逻辑，就可能分别延迟开启或延迟关闭。因此一致前馈并非“没有适应所以无趣”，它回答的是另一个时间功能问题。

\subsection{汤超参与的 Cell 工作：从“所有图”转向“目标功能”}
Ma、Trusina、El-Samad、Lim 与汤超的 \emph{Defining Network Topologies that Can Achieve Biochemical Adaptation}（Cell, 2009，参考文献~6）提供了一种研究思路：先定功能标准，再在统一动力学规则下枚举。研究节点是总量固定、在活性与非活性状态之间转换的酶；连边按 Michaelis--Menten 反应建模，和本文的转录产生--降解模型不同。

三节点的九条潜在有向边给出 $3^9$ 种符号图。当输入固定在 $A$、输出固定在 $C$ 时，没有因果路径的图必须缺少 $A\to C$，而且 $A\to B$、$B\to C$ 不能同时存在。这两条边共有 $5$ 种不形成路径的组合，其余六条边任取，故筛去 $5\times3^6$ 种，剩下 $3^9-5\times3^6=16\,038$ 种。论文对这些拓扑各采样 $10\,000$ 组参数，用灵敏度与适应精度筛选。每图采样中超过 $10$ 组参数达标的 $395$ 个较稳健拓扑均含缓冲节点负反馈（NFBLB）或比例节点不一致前馈（IFFLP）核心。这里的数目是该模型与筛选规则的结果，不是所有三基因系统的功能计数。

为了理解筛选到底在看什么，设输入从 $I_1$ 变为 $I_2$，输出从初始稳态 $O_1$ 出发，达到最明显偏离的位置 $O_{\rm peak}$，最后趋于 $O_2$。取非零的基线和输入变化，定义
\[
\Delta_I=\frac{|I_2-I_1|}{|I_1|},\qquad
S=\frac{|O_{\rm peak}-O_1|/|O_1|}{\Delta_I},\qquad
E=\frac{|O_2-O_1|/|O_1|}{\Delta_I},\qquad P=E^{-1}.
\]
$S$ 是\emph{灵敏度}，问输出有没有明显响应；$P$ 是\emph{适应精度}，问最终偏差是否很小。论文采用 $S>1,P>10$ 的功能门槛。精确适应对应 $E=0$，形式上 $P$ 无限大，但这不是数值采样直接证明的性质。

一个始终不动的输出可以有很小的 $E$，却不是良好的变化探测器；一个不断振荡的输出也不能简单用最终平台 $O_2$ 评价。论文补充材料排除了持续振荡，并对衰减摆动设置了额外规则。对“适应”的筛选失败不等于对“振荡”的研究失败：在前面振荡例子中有用的时钟，在这里恰恰可能被排除。

\subsection{缓冲节点负反馈：把偏差积累起来，直到输出恢复}
为什么普通的负反馈不自动给出精确适应？因为一个通常的恢复变量会同时遗忘过去的偏差：产生与降解平衡时，允许输出停在一个小的残余误差上。要把这个误差消掉，调节变量需要在误差存在时继续改变。

下面用本文的简化三变量模型说明 NFBLB 的这一作用。$A$ 接收输入，$C$ 是输出，$B$ 是调节变量；$C$ 激活 $B$，$B$ 抑制 $C$。令
\begin{align}
\dot A&=\rho\{u(1-A)-A\},\\
\dot B&=\kappa(C-C_0),\\
\dot C&=k\{A(1-C)-BC\},\qquad \rho,\kappa,k>0.\label{eq:buffer_demo}
\end{align}
$A,C$ 可以理解为活性比例；$B$ 是缓冲变量，在本例的工作范围内也保持正且小于 $1$。第一式使 $A^*=u/(1+u)$。关键在第二式：当 $C>C_0$，$B$ 继续积累并压低 $C$；当 $C<C_0$，$B$ 下降，抑制减弱；只有 $C=C_0$ 时，$B$ 才停止调整。因此
\[
C^*=C_0,\qquad B^*=\frac{1-C_0}{C_0}A^*.
\]
输入变大后，长期改变的是 $A^*$ 与 $B^*$，输出仍为 $C_0$。这不是两个支路恰巧抵消一次，而是反馈持续校正误差。

把第二式积分，得到
\[
B(t)=B(0)+\kappa\int_0^t[C(s)-C_0]\dd s.
\]
所以 $B$ 保存的是过去输出误差的累积，称为\emph{积分反馈}。这里的“缓冲”不只是响应慢，而是 $\dot B$ 在工作近似下不直接依赖 $B$ 自己。饱和的酶反应能提供近似恒定的转化速率，从而产生这类条件；它不会无条件地出现在所有负反馈回路中。式~\eqref{eq:buffer_demo}是机制模型，不是论文全套速率方程；若 $B$ 达到上下限，线性积分近似就失效，必须回到完整的转化反应。

这个例子的恢复也能算出来。因为 $A$ 不受下游影响，线性化后它的一根特征值为 $-\rho(1+u)$，而 $(B,C)$ 部分为
\[
J_{BC}=\begin{pmatrix}0&\kappa\\-kC_0&-k(A^*+B^*)\end{pmatrix}.
\]
迹为负，行列式 $\kappa kC_0>0$，故只要 $0<C_0<1$ 且 $B^*$ 在缓冲近似的内部工作范围，稳态就局部渐近稳定。输出基线固定与基线附近能恢复，是这里分别检查过的两件事。

若给缓冲变量加入“泄漏”，改成
\[
\dot B=\kappa(C-C_0)-\varepsilon_B(B-B_{\rm ref}),\qquad\varepsilon_B>0,
\]
稳态就允许 $C^*-C_0=(\varepsilon_B/\kappa)(B^*-B_{\rm ref})$。它仍是负反馈，也仍可能有瞬态峰，却通常留下随输入变化的基线偏差。缓冲变量慢并不足够，是否能持续积累误差才是区别。

\subsection{比例节点前馈：长期随输入增加，恰好抵消直接支路}
前馈的思路不同：没有输出回传的校正，而是让一条支路产生与输入成比例的反向作用。沿用同样的输入节点 $A$，构造另一个简化模型：
\begin{align}
\dot A&=\rho\{u(1-A)-A\},\\
\dot B&=k_BA-\gamma_BB,\\
\dot C&=k_AA(1-C)-k_CB C,\qquad k_B,\gamma_B,k_A,k_C>0.
\end{align}
$A$ 直接激活 $C$，又通过 $B$ 抑制 $C$。令 $\eta=k_B/\gamma_B$，稳态下
\[
B^*=\eta A^*,\qquad
0=A^*\{k_A(1-C^*)-k_C\eta C^*\}.
\]
在 $A^*>0$ 的工作范围内，公共因子 $A^*$ 可以约去，得到
\[
C^*=\frac{k_A}{k_A+k_C\eta}.
\]
这就是比例节点的作用：输入增加多少，抑制支路长期也按固定比例增加多少，两条支路在输出方程中共享输入因子。$B$ 较慢时，直接支路先起作用，随后比例支路跟上，于是输出先偏离、再回到基线。

该模型是逐级驱动的：$A$ 先稳定，随后 $B$ 稳定，最后 $C$ 的恢复率为 $k_AA^*+k_CB^*>0$，所以在给定正输入及上述工作范围内三者都收敛。若把 $B$ 的产生换成 $k_BA/(K+A)$，则 $B^*/A^*=k_B/[\gamma_B(K+A^*)]$ 不再恒定，输出基线也通常随输入改变。可见不一致前馈的边符号还要配上合适的支路比例关系。

\begin{figure}[H]\centering
\begin{tikzpicture}
\node[gene] (a) at (0,1.3) {$A$};\node[gene] (c) at (3,1.3) {$C$};\node[gene] (b) at (3,-.4) {$B$};
\draw[act] (a)--(c);\draw[act] (c) to[bend left=22] (b);\draw[rep] (b) to[bend left=22] (c);
\node at (1.5,-1.15) {缓冲：$\dot B\propto C-C_0$};
\node[gene] (aa) at (6,1.3) {$A$};\node[gene] (cc) at (9,1.3) {$C$};\node[gene] (bb) at (7.5,-.4) {$B$};
\draw[act] (aa)--(cc);\draw[act] (aa)--(bb);\draw[rep] (bb)--(cc);
\node at (7.5,-1.15) {比例：$B^*\propto A^*$};
\end{tikzpicture}
\caption{两种适应机制的节点角色。左边由输出误差驱动缓冲变量持续校正；右边由输入驱动比例变量，长期抵消直接输入作用。两者都允许 $B$ 的稳态随输入改变。}
\end{figure}
\begin{figure}[H]\centering\includegraphics[width=\linewidth]{figures/adaptation_mechanisms.pdf}
\caption{本文简化模型的数值比较。左：缓冲反馈与比例前馈都先响应再恢复，带泄漏的缓冲反馈改变最终基线。中：调节变量 $B$ 保留输入变化。右：理想机制的稳态输出与输入无关，泄漏模型则不然。曲线按各自初始基线归一化；并非复现论文的全网参数扫描。}\label{fig:adaptmech}
\end{figure}
本图的示例参数为 $\rho=3$、输入 $u:0.5\to0.8$。缓冲模型取 $C_0=0.4,\kappa=0.3,k=2$，泄漏模型另取 $\varepsilon_B=0.08$；比例模型取 $k_B=\gamma_B=0.25,k_A=k_C=2$。这些参数用于展示基线恢复与泄漏反例，没有按论文的 $S>1$ 灵敏度门槛选取，不能把本图直接计入作者的功能筛选结果。完整方程与检查见 \texttt{code/extended\_examples.py}。

\subsection{为什么功能枚举不能只看节点数与连边？}
论文的补充材料进一步显示：同一个负反馈环，把哪个节点当作输出会改变是否适应；带正自边的特定速率条件还可让 $\log B$ 累积输出误差。例如补充材料的一种近似为 $\dot B=\nu B(C_0-C)$，在 $B>0$ 时可写成
\[
\frac{\dd\log B}{\dd t}=\nu(C_0-C).
\]
于是存储累计误差的是 $\log B$，内部稳态仍由 $C=C_0$ 决定。这个例子涉及饱和速率条件及随调控变量改变的损失作用，不能直接套用固定一阶降解模型的分类结论。因此“缓冲节点”与“比例节点”是动态角色，不只是绘图中的某个位置。两个例子也提醒我们，正自反馈既可能存储状态，也可能参与适应；功能要由方程中的作用方式来区分（参考文献~7）。

论文还扫描了固定输入节点 $A$、输出节点 $C$ 的 $81$ 种双节点酶网络，在所扫参数范围内没有得到符合其标准的适应。它没有证明任意两变量系统都不能适应：比如把输入直接施加于输出，再用另一个变量积分误差，就改变了输入、输出角色；换用不同降解机制也改变了模型族。本文的两节点分类讨论的是恢复、记忆和 Hopf 入口，而论文在同样数量的结构上问的是另一个有指定标准的功能问题。

“稳健”也要和“精确”区分。精确适应问某个模型的稳态输出是否在输入区间内恒定；稳健性问参数不必调到孤立的特殊数值时，这个功能是否仍能实现。数值上可以记录同一拓扑的 $N$ 组参数中有多少组符合标准，写成 $Q$ 或比例 $Q/N$。$Q$ 依赖参数范围、采样规则、输入步骤和功能门槛，不能当作拓扑唯一的内在属性。它也不要求输出基线对所有参数变化都不变：同一组参数下适应不同输入，与换参数后保持同一个基线，是两个不同要求。

所以更彻底的三节点研究应同时记录\emph{拓扑、节点角色、速率机制、参数、功能标准}。一致前馈可以筛选信号持续时间，不一致前馈可以形成脉冲并在比例条件下适应，负反馈既可能恢复或振荡，也可能在缓冲条件下适应。对一项功能不合格的网络，可能恰好实现另一项功能；这比给所有图固定贴上“有趣”或“无趣”的标签更有解释力。

\section{分子数较少时：从浓度轨迹转向概率分布}
前面的轨迹把表达量当作连续浓度。在分子数较多、涨落相对较小时，这能突出反馈的主要作用；但同一个产生--降解机制在低分子数下也会不断随机变化。此时，长期“稳定”不再意味着每条轨迹停在同一个数值，而是分子数的概率分布保持不变。先回到最简单的无反馈反应，看看这一转换怎样进行。

\subsection{同一组反应怎样产生主方程？}
考虑 $k>0,\gamma>0$ 的出生--死亡过程
\begin{equation}\varnothing\xrightarrow{k}X,\qquad X\xrightarrow{\gamma}\varnothing.\end{equation}
令 $P_n(t)$ 表示时刻 $t$ 恰有 $n$ 个分子的概率。在短时间 $\dd t$ 内，出生使 $n-1\to n$，跃迁率为 $k$；死亡使 $n+1\to n$，跃迁率为 $\gamma(n+1)$；从 $n$ 离开的总率为 $k+\gamma n$。流入减去流出，只保留 $\dd t$ 的一阶项，就得到
\begin{equation}\dot P_n=kP_{n-1}+\gamma(n+1)P_{n+1}-(k+\gamma n)P_n,\qquad P_{-1}=0.\end{equation}
零边界没有向负分子数的跃迁。确定性方程描述一个浓度的导数，而主方程同时描述所有可能分子数之间的概率转移。

\subsection{稳态为什么是泊松分布？}
先寻找长期不变的分布。定义从 $n$ 到 $n+1$ 的净概率流
$j_n=kP_n-\gamma(n+1)P_{n+1}$，主方程就是 $\dot P_n=j_{n-1}-j_n$。在稳态各相邻概率流必须相等，而零边界给出 $j_{-1}=0$，所以所有 $j_n=0$。这就推出相邻态的\emph{详细平衡}：每一对相邻分子数之间，向上与向下转移的概率流分别相等。它比只要求总分布不再变化更具体。递推得到
\[
P_n=P_0\frac{\mu^n}{n!},\qquad\mu=\frac k\gamma.
\]
归一化要求 $1=P_0\sum_{n=0}^{\infty}\mu^n/n!=P_0e^\mu$，从而
\begin{equation}P_n=e^{-\mu}\frac{\mu^n}{n!}.\end{equation}
因此泊松分布是这条链的唯一归一化稳态；平均分子数 $\mu$ 与此前产生--降解的平衡值对应，但单次观测仍可偏离平均值。
\begin{figure}[H]\centering\includegraphics[width=.74\linewidth]{figures/birth_death.pdf}\caption{基本出生--死亡模型的稳态泊松分布。平均分子数越低，涨落相对于平均值越明显。}\end{figure}

\subsection{平均值、涨落幅度与恢复时间怎样联系？}
为了比较主方程与浓度模型，可以直接计算平均值与方差。先假设平均值与方差都是有限的，数学上称为\emph{有限二阶矩}：$\langle n^2\rangle<\infty$，所以分子数平方的平均值不会无限大。下面分别关心分子数 $h(n)=n$ 和它的平方 $h(n)=n^2$。把每次出生、死亡对这些数量的改变，按照发生概率加起来，就有
\[
\frac{\dd\langle h(n)\rangle}{\dd t}
=\left\langle k[h(n+1)-h(n)]+\gamma n[h(n-1)-h(n)]\right\rangle.
\]
先取 $h=n$，得到
\begin{equation}\frac{\dd\langle n\rangle}{\dd t}=k-\gamma\langle n\rangle.\end{equation}
这个平均值方程与确定性产生--降解模型具有相同形式，是线性反应下的精确关系。再取 $h=n^2$，得到
\[
\frac{\dd\langle n^2\rangle}{\dd t}=2k\langle n\rangle+k-2\gamma\langle n^2\rangle+\gamma\langle n\rangle.
\]
减去 $\dd\langle n\rangle^2/\dd t$，就有
\begin{equation}\frac{\dd\operatorname{Var}(n)}{\dd t}=k+\gamma\langle n\rangle-2\gamma\operatorname{Var}(n).\end{equation}
在稳态代入 $\langle n\rangle=\mu=k/\gamma$，得到 $\operatorname{Var}(n)=\mu$。因此变异系数为 $1/\sqrt\mu$，分子数越低，相对内禀涨落越大。

涨落还具有时间结构。由平均值方程，从给定 $n(0)$ 出发的条件均值为
\[
\mathbb E[n(t)\mid n(0)]=\mu+(n(0)-\mu)e^{-\gamma t}.
\]
减去 $\mu$，再乘以 $n(0)-\mu$ 并在稳态平均，得到
\begin{equation}\operatorname{Cov}(n(t),n(0))=\mu e^{-\gamma t},\qquad t\geq0.\end{equation}
所以同一个 $1/\gamma$ 既描述平均浓度的恢复，也描述这条线性反应链的涨落相关时间。若产生率带有非线性反馈，通常有 $\langle F(n)\rangle\ne F(\langle n\rangle)$，平均值不再自动满足原来的非线性浓度方程。

\subsection{两个概率峰是否意味着两个确定性稳定状态？}
前面的自激活和互抑制开关有两个确定性吸引态，但在随机模型中，两个概率峰还可能来自另一种时间结构：启动子在两种产生率之间缓慢切换。取关闭与开启状态的产生率为 $k_0,k_1$，共享降解率 $\gamma$，切换率为 $k_{\rm on},k_{\rm off}$。

如果两种切换均远慢于 $\gamma$，蛋白质常有足够时间接近当前启动子状态的条件平衡。二态链的稳态流量平衡给出开启概率 $q=k_{\rm on}/(k_{\rm on}+k_{\rm off})$，整体分布的领先阶近似于是为
\[
P_n\simeq(1-q)\operatorname{Pois}_{k_0/\gamma}(n)+q\operatorname{Pois}_{k_1/\gamma}(n).
\]
若两个均值相距足够远且两权重均显著，这个混合分布可以双峰；每个固定启动子状态下的浓度模型却都只有唯一稳定的稳态。该式是慢切换近似，有限切换速率下不是精确公式。由此可见，解释两个峰时还需要看时间序列、驻留时间和扰动响应，判断状态选择来自确定性反馈，还是来自慢切换或其他机制。

\section{回到具体生物学：乳糖操纵子中的有效反馈}
至此，最小模型已经解释了记忆、振荡、适应与随机涨落各自可能来自什么机制。把它们用于真实系统时，第一步是识别生物过程怎样对应到模型变量和调控函数。乳糖操纵子提供了一个例子：反馈并不只发生在基因之间，还经过摄取与代谢。
大肠杆菌 lac 系统包括 LacI 阻遏蛋白、lacZYA 转录单位以及乳糖/异乳糖相关代谢与摄取过程。lacI 具有自身的转录调控，不属于 lacZYA 同一多顺反子转录单位。异乳糖降低 LacI 的阻遏活性；CAP--cAMP 及葡萄糖相关效应调节转录和摄取。

在最简单的结合占据近似下，活性 LacI 浓度为 $R_a$ 时，解除阻遏的分数可写作 $K_d/(K_d+R_a)$，从而给出有效转录率 $F(R_a)$。进一步引入 LacY 蛋白 $y$ 与细胞内诱导物 $a$，
\begin{equation}\dot y=F(a)-\gamma_y y,\qquad\dot a=v_{\rm in}(y,L_{\rm ext})-v_{\rm out/met}(a,y).\end{equation}
这体现了 ``LacY 增强摄取 $\to$ 诱导物增加 $\to$ 阻遏解除 $\to$ LacY 增加'' 的有效正反馈，但真正的乳糖代谢、透过酶动力学、葡萄糖作用和生长稀释会改变反馈强度。特别应区分代谢性乳糖与非代谢型诱导物的实验情形，不能笼统声称 lac 网络必然双稳态。

\begin{figure}[H]\centering
\begin{tikzpicture}
\node[gene] (y) at (0,0) {LacY};
\node[gene] (a) at (3,0) {胞内诱导物};
\node[gene] (r) at (6,0) {活性 LacI};
\node[gene] (f) at (9,0) {lacZYA 表达};
\draw[act] (y)--node[above]{摄取}(a);
\draw[rep] (a)--(r);\draw[rep] (r)--(f);
\draw[act] (f.south) -- ++(0,-.7) -| (y.south);
\end{tikzpicture}
\caption{lac 系统的简化有效反馈关系。两次抑制形成解除阻遏，闭环总体为正；该图省略了代谢、葡萄糖作用和稀释，不据此断言双稳态。}
\end{figure}

为了把这条生物学路径与前面的符号分析接上，可以逐步计算导数。单个位点近似下 $F(R_a)=r_{\rm leak}+(r_{\rm max}-r_{\rm leak})K_d/(K_d+R_a)$，当 $r_{\rm max}>r_{\rm leak}$ 时 $\partial F/\partial R_a<0$。
若有效活性阻遏蛋白 $R_a(a)$ 随诱导物下降，则链式法则给出 $\dd F/\dd a=(\partial F/\partial R_a)(\dd R_a/\dd a)>0$。
再结合 $\partial v_{\rm in}/\partial y>0$，即可识别摄取支路的正反馈。完整第二个方程的交叉导数却是
$\partial\dot a/\partial y=\partial_yv_{\rm in}-\partial_yv_{\rm out/met}$，其符号未必正；这正是代谢性诱导物与仅摄取的理想模型之间需要检查的环节。

这个例子把最小回路的用法具体化了：先识别产生率与中间变量，再检查有效反馈的符号和强度，最后分析对应参数下的动力学。接下来还要问，多个因子共同作用于启动子时怎样合并？局部回路受到其余网络影响时，独立分析还能保留多少结论？

\section{从局部回路到真实网络：还需要补上什么？}
lac 的摄取支路已经说明，简化图上的一条边可能压缩了多个过程。真实网络还存在组合输入、与其余变量的动态耦合，以及统计意义上的模体富集。理解这些层面，才能判断最小模型的解释在何种条件下可以迁移。
\subsection{组合调控：符号拓扑不等于启动子逻辑}
两个转录因子 $X,Y$ 结合启动子 $D$ 的四个状态为 $D,DX,DY,DXY$。在特定准平衡模型中，统计权重为
\begin{equation}1,\quad x/K_x,\quad y/K_y,\quad\omega xy/(K_xK_y),\end{equation}
$\omega$ 表示协同性。令 $Z=1+x/K_x+y/K_y+\omega xy/(K_xK_y)$，则
\begin{equation}F(x,y)=\frac{r_0+r_xx/K_x+r_yy/K_y+r_{xy}\omega xy/(K_xK_y)}{Z}.\end{equation}
同样两条激活边，由 $r_s$ 的不同可形成 AND-like、OR-like 等不同调控逻辑。仅有网络连边不能确定动力学方程。

具体的差别可以从状态概率直接算出来。写 $a=x/K_x,b=y/K_y$。准平衡时四种状态的概率是各自统计权重除以总权重 $Z$，所以 $F=\sum_s r_s p_s$，即前述有理式。$\omega=1$ 时 $Z=(1+a)(1+b)$，记 $p_x=a/(1+a),p_y=b/(1+b)$。
若只有双结合态转录，$r_{xy}=r$、其余为零，得到
$F_{\rm AND}=rp_xp_y$；若除空态以外均以速率 $r$ 转录，得到
$F_{\rm OR}=r(p_x+p_y-p_xp_y)$。
二者对 $x,y$ 均非减，却有不同的低输入响应。极限 $a,b\ll1$ 或 $\gg1$ 将连续调控近似为布尔真值表。若 $r_{xy}$ 小于单结合态速率，增加另一个因子甚至可能降低输出；因此指定某些 $r_s$ 后还要实际检查偏导数，不能只凭结合就称为激活。
\begin{figure}[H]\centering
\begin{tikzpicture}[node distance=20mm]
\node[gene] (d) {$D$};\node[gene,right=of d] (dx) {$DX$};
\node[gene,below=of d] (dy) {$DY$};\node[gene,right=of dy] (dxy) {$DXY$};
\draw[<->] (d)--node[above]{$X$}(dx);\draw[<->] (d)--node[left]{$Y$}(dy);
\draw[<->] (dx)--node[right]{$Y$}(dxy);\draw[<->] (dy)--node[below]{$X$}(dxy);
\node[align=left,right=15mm of dx,anchor=west] {状态权重：$1,a,b,\omega ab$\\状态转录率：$r_0,r_x,r_y,r_{xy}$};
\end{tikzpicture}
\caption{两种调控因子的启动子占据状态。图示可逆结合；协同性必须与相容的结合常数共同指定，最终输出由状态概率与各状态转录率共同决定。}
\end{figure}

\subsection{模块嵌入：局部稳定为何不能直接拼装？}
将已分析的局部变量记为 $\bm x$，其余网络记为 $\bm z$：
\begin{equation}\dot{\bm x}=\bm f(\bm x,\bm z),\qquad\dot{\bm z}=\bm g(\bm x,\bm z).\end{equation}
在稳态附近，两个模块的小偏差怎样相互影响，可以整理成四块： $J=\begin{psmallmatrix}A&B\\C&D\end{psmallmatrix}$。其中 $A,D$ 描述两个模块各自的变化，$B,C$ 描述它们如何影响对方。即使两个模块单独运行都能恢复，连接起来也可能让某些偏差不断放大，因此不能只用单独模块的结果判断整个网络。只有在外部输入近似固定、存在明确时间尺度分离或耦合可控时，局部回路才可能构成近似独立运行的模块。这解释了为何最小回路有用，但不能无条件拼装成真实系统。

一个两维线性扰动例子足以看出耦合的影响。取 $A=D=-1$、$B=C=2$。单独运行时，两块的偏差都按 $e^{-t}$ 衰减。连接后，描述小偏差的雅可比矩阵
$\begin{psmallmatrix}-1&2\\2&-1\end{psmallmatrix}$ 的特征值却是 $1,-3$，所以耦合足以使系统失稳。这是线性扰动系统的反例，不把可取负值的扰动直接解释为绝对浓度。
要继续计算连接后的影响，可以先从外部模块的方程中解出它的偏差，再代回局部模块。只要 $\lambda I-D$ 有逆矩阵，这一代入就给出
\[
\det(\lambda I-J)=\det(\lambda I-D)
\det\left[\lambda I-A-B(\lambda I-D)^{-1}C\right].
\]
这个消去外部变量后的组合称为\emph{Schur 补}：它把外部变量的响应折算成局部方程中的修正。这里的 $B(\lambda I-D)^{-1}C$ 就是“局部影响外部，再由外部传回局部”的整条反馈路径。代入的矩阵计算相当于左乘 $\begin{psmallmatrix}I&B(\lambda I-D)^{-1}\\0&I\end{psmallmatrix}$，把右上块消成零；这个操作不改变行列式，剩下的两块分别取行列式再相乘，就得到上式。这个式子说明：外部模块传回来的影响取决于偏差变化的快慢，即时间因子 $e^{\lambda t}$ 中的 $\lambda$；不能简单地当成一个固定输入。
如果外部变量恢复得远快于局部变量，就可以近似认为它们总能及时跟上。令外部变量的变化率接近零，即 $0=C\delta\bm x+D\delta\bm z$ 给出 $\delta\bm z=-D^{-1}C\delta\bm x$，从而 $A_{\rm eff}=A-BD^{-1}C$。所以即使把外部变量从方程中省去，它们的影响仍留在 $-BD^{-1}C$ 这一项里，局部模块的实际恢复行为通常也会改变。

\subsection{网络模体：结构富集怎样与功能分析连接？}
明确了调控函数与网络背景之后，再讨论一个小回路为何在真实网络中反复出现。
真实网络中常统计自调控、前馈回路、反馈环等小拓扑是否相对某一合理随机网络零模型富集。富集不等于因果功能，也不能自动认定局部回路独立于网络其余部分。较妥当的问题是：哪些局部动力学性质在给定输入变化和网络背景下依然可观察？

\section{沿着主线回看：每次增加了什么机制？}
从产生--降解出发，单节点正反馈通过多个稳态提供记忆，双节点互抑制通过放大表达差异形成开关，三节点负反馈增加动态环节而能使带摆动的小偏差从衰减转为增长，快慢前馈则以支路竞争形成瞬态响应。转向分子计数后，同一个长期平均值还对应一整个涨落分布；回到真实网络时，组合逻辑与模块耦合进一步决定这些局部机制能否被观察到。下表把各步得到的结论与它们的适用边界放在一起。
\begin{center}\begin{tabular}{>{\raggedright\arraybackslash}p{3.1cm}>{\raggedright\arraybackslash}p{6cm}>{\raggedright\arraybackslash}p{5.6cm}}\toprule
模型 & 由推导得到的行为 & 迁移结论时需注意\\\midrule
线性产生--降解 & 唯一指数稳定的稳态；出生--死亡离散模型为泊松稳态 & 泊松结论依赖恒定出生与一阶死亡\\
一维单调自抑制 & 唯一稳定的稳态；不能自行维持振荡 & 输入固定且没有延迟时成立\\
一维自激活 & 足够非线性时可能双稳态 & 多交点还需足够非线性与合适参数\\
双组分互相调控 & 正/负回路符号影响特征值和稳定性 & 稳定性仍依赖增益及具体产生率\\
三节点负反馈 & 小偏差可从衰减摆动转为增长摆动 & 稳定周期轨道需要非线性分析\\
前馈回路 & 可实现脉冲、滤波和特定条件下适应 & 精确适应还需稳态支路增益抵消\\
\bottomrule\end{tabular}\end{center}

分析一个新回路时，也可以沿这条路线：先从反应规则建立产生率，说明哪些快变量被省去；再寻找稳态，检查偏差会恢复还是放大；最后考虑分子涨落与其余网络的影响，判断模型能解释哪些实验现象。

\clearpage
\section*{附录：81 种两节点拓扑逐图索引}\addcontentsline{toc}{section}{附录：81 种两节点拓扑逐图索引}
四条边的编号次序固定为 $(X\to X,Y\to Y,X\to Y,Y\to X)$。以下三张表展示全部 $81$ 张带标签图，每张只出现一次。R 编号表示结构组，不表示某组参数已经实现的行为。枚举脚本检查了 $81$ 个唯一图、$45$ 个节点交换类别及正文八组的计数；拓扑信息可由 \texttt{code/topology\_inventory.json} 直接复用。
\input{figures/topology_atlas_1.tex}
\input{figures/topology_atlas_2.tex}
\input{figures/topology_atlas_3.tex}

\clearpage\section*{参考文献与延伸阅读}\addcontentsline{toc}{section}{参考文献与延伸阅读}
\begin{enumerate}[leftmargin=*,itemsep=2pt,parsep=0pt]\small
\item Bintu, L. et al. (2005). Transcriptional regulation by the numbers: models. \emph{Current Opinion in Genetics \& Development}, 15, 116--124. DOI: \href{https://doi.org/10.1016/j.gde.2005.02.007}{\nolinkurl{10.1016/j.gde.2005.02.007}}.
\item Gardner, T. S., Cantor, C. R., \& Collins, J. J. (2000). Construction of a genetic toggle switch in Escherichia coli. \emph{Nature}, 403, 339--342. DOI: \href{https://doi.org/10.1038/35002131}{\nolinkurl{10.1038/35002131}}.
\item Elowitz, M. B. \& Leibler, S. (2000). A synthetic oscillatory network of transcriptional regulators. \emph{Nature}, 403, 335--338. DOI: \href{https://doi.org/10.1038/35002125}{\nolinkurl{10.1038/35002125}}.
\item Shen-Orr, S. S. et al. (2002). Network motifs in the transcriptional regulation network of Escherichia coli. \emph{Nature Genetics}, 31, 64--68. DOI: \href{https://doi.org/10.1038/ng881}{\nolinkurl{10.1038/ng881}}.
\item Alon, U. (2007). Network motifs: theory and experimental approaches. \emph{Nature Reviews Genetics}, 8, 450--461. DOI: \href{https://doi.org/10.1038/nrg2102}{\nolinkurl{10.1038/nrg2102}}.
\item Ma, W., Trusina, A., El-Samad, H., Lim, W. A., \& Tang, C. (2009). Defining network topologies that can achieve biochemical adaptation. \emph{Cell}, 138, 760--773. DOI: \href{https://doi.org/10.1016/j.cell.2009.06.013}{\nolinkurl{10.1016/j.cell.2009.06.013}}.
\item Ma, W. et al. (2009). Supplemental Data for \emph{Defining Network Topologies that Can Achieve Biochemical Adaptation}. 重点参看缓冲节点、输入/输出位置与对数积分的讨论。\href{https://ucsf-cmp.github.io/lim-lab/papers/pdfs/wm_2009_supp.pdf}{作者补充材料}。
\item Tsai, T. Y.-C., Choi, Y. S., Ma, W., Pomerening, J. R., Tang, C., \& Ferrell, J. E., Jr. (2008). Robust, tunable biological oscillations from interlinked positive and negative feedback loops. \emph{Science}, 321, 126--129. DOI: \href{https://doi.org/10.1126/science.1156951}{\nolinkurl{10.1126/science.1156951}}.
\end{enumerate}
\end{document}
